\subsection{Actual ruled contractions across triple points}

\begin{statement}{C.2.UN.74}{Proved}{A section and ruling fibres give the exact transverse curve blowup}
Let $f:U\to\Delta$ be a proper degeneration of connected
complex two-tori from a smooth threefold, with SNC reduced
central support and arbitrary positive full multiplicities.
Let $T$ be a component with a specified holomorphic
projective-line bundle $p:T\to B$ over a smooth compact
curve. Suppose its boundary in the reduced support is
exactly
\[
           C+F_1+\cdots+F_r,\qquad
           C\ \text{a section},\quad F_j=p^{-1}(b_j),
\tag{C.RC.1}
\]
where the $b_j$ are distinct. The section is the
intersection with one component $P$, and $F_j$ is the
intersection with a component $Q_j$. The components
$Q_j$ may repeat at different $b_j$, but none is $P$.
All these are actual reduced transverse intersection
divisors; their intersections $C\cap F_j$ are the
specified triple points. Let the full multiplicities
of $T,P,Q_j$ be respectively $m_T,m_P,n_j$.

There is an ordinary smooth curve contraction along
the ruling, inverse to an ordinary blowup and equal
to an isomorphism off $T$, exactly when
$m_T=m_P$. In particular this always holds for a
full reduced fibre. Set their common value equal
to $m$. The resulting actual map is
\[
              \pi:U\longrightarrow U',\qquad
              f=f'\pi,\qquad \pi(T)=B\subset P'.
\tag{C.RC.2}
\]
The threefold $U'$ is smooth and $f'$ is proper with
the same punctured torus family. Its reduced central
support is SNC. The map $P\to P'$ is an isomorphism.
Each $Q_j\to Q'_j$ is the ordinary point blowup at
$b_j\in B$, with all indicated points blown up when
the same component repeats. The curve $B$ lies in
the smooth surface $P'$, avoids every triple point
of the descended support, and meets its double
curves transversely, exactly at the $b_j$.

All full multiplicities are retained: the parent
has multiplicity $m$ and $Q'_j$ has multiplicity
$n_j$. At $b_j$ there are actual coordinates
$(p_0,q_0,r_0)$ with
\[
 \begin{gathered}
 P'=(p_0=0),\quad Q'_j=(q_0=0),\quad
 B=(p_0=r_0=0),\qquad f'=p_0^m q_0^{\,n_j},\\
 p_0=xy,\quad r_0=x,\quad q_0=z,\qquad
                         f=x^m y^m z^{\,n_j}
 \end{gathered}
\tag{C.RC.3}
\]
in the blowup chart containing $C\cap F_j$.
Every other component and every neighbourhood away
from the indicated exceptional configurations are
unchanged.

Write
$L=N_{B/P'}$ and $M=N_{P'/U'}|_B$.
The actual normal extension and full-fibre
trivialization are
\[
 \begin{gathered}
 0\longrightarrow L\longrightarrow
       V=N_{B/U'}\longrightarrow M\longrightarrow0,
 \qquad (T,C)=(\PP_B(V),\PP_B(L)),\\
 N_{T/U}\simeq\OO_{\PP_B(V)}(-1),\qquad
 M^{\otimes m}\simeq
             \OO_B\left(-\sum_{j=1}^r n_jb_j\right).
 \end{gathered}
\tag{C.RC.4}
\]
Thus the full fibre weights on the ruling-fibre
boundary enter the actual root line $M$; they are
not discarded. In particular $m$ divides
$\sum_jn_j$. For actual canonical orders $k_P,k_{Q_j}$
on the descended family, the section subline is
\[
 L\simeq K_B\otimes M^{-(k_P+1)}
                    \otimes\OO_B\left(-\sum_jk_{Q_j}b_j\right),
             \qquad k_T-k_P=1.
\tag{C.RC.5}
\]
The entire extension in (C.RC.4), with its endpoint
identifications, is retained.

For a full reduced fibre, put $Z=\sum_jb_j$. Then
these maps specialize to
\[
 \begin{gathered}
 M\simeq\OO_B(-Z),\qquad
 L\simeq K_B\otimes
              \OO_B\left(\sum_j(k_P-k_{Q_j}+1)b_j\right),\\
 df'|_V:\ V\twoheadrightarrow
            \OO_B(-Z)\hookrightarrow\OO_B.
 \end{gathered}
\tag{C.RC.6}
\]
If in addition $T$ is a local maximum of the
canonical order and $r>0$, the extension splits:
$V\simeq L\oplus M$, with its specified subline.

For the underlying reduced supports,
\[
              e(S)=e(S')+2-2g(B),\qquad
              t(S)=t(S')+r,
\tag{C.RC.7}
\]
where $t$ counts triple points. The same contraction
exists on a compact total space. Conversely an
ordinary blowup of a smooth curve in a single
central component, transverse to its double curves
and avoiding the triple locus, has exactly the
maps, markings and multiplicities displayed above.
\end{statement}

\begin{proof}
\par\medskip\noindent\textbf{The actual normal degree and the corrected contraction theorem.}
Restrict the principal full central fibre to $T$.
For a ruling fibre $\ell$ away from the marked
values, the resulting line isomorphism gives
\[
       m_T\deg(N_{T/U}|_\ell)=-m_P,\qquad
       N_{T/U}^{\,m_T}
           \simeq\OO_T\left(-m_PC-\sum_jn_jF_j\right).
\tag{C.RC.8}
\]
If $m_T=m_P$, the restriction is precisely
$\OO_{\PP^1}(-1)$. The corrected theorem of
\cite[\S2, Theorem~2' and Proposition~9']{FujikiNakano1971},
in the form used in UN64, supplies the actual
smooth blowdown along $p:T\to B$. It identifies
$U$ with the ordinary blowup of a smooth curve
in a smooth threefold $U'$. Conversely an
ordinary such blowup has normal degree minus
one, so (C.RC.8) forces $m_T=m_P$.

The equality $\pi_*\OO_U=\OO_{U'}$ gives the
unique holomorphic $f'$ in (C.RC.2). Its
properness follows because the inverse image
of a compact set in the disc is the image
under $\pi$ of the corresponding compact
set for $f$. The smooth fibres are unchanged.
It remains to prove the claimed structure
at the centre rather than infer it from a
numerical normal degree.

\par\medskip\noindent\textbf{The parent image is smooth and its whole surface is retained.}
The proper image $P'=\pi(P)$ is a hypersurface
in the smooth target. Its actual Cartier
pullback has the form
\[
                            \pi^*P'=P+aT.
\tag{C.RC.9}
\]
Restriction to an ordinary ruling fibre
gives $0=1-a$, since $P\cap T=C$ is a
section and $T|_\ell=\OO_\ell(-1)$.
Thus $a=1$.

A local defining section of $P'$ vanishes
to first order along the generic centre.
Its first normal term on $\PP_B(N_{B/U'})$
has zero divisor exactly $P\cap T=C$,
by (C.RC.9). This term cannot vanish at
any whole fibre: $C$ is a section and
contains none. Its differential in the
normal directions is therefore nonzero
at every point of $B$. Hence $P'$ is
smooth along $B$. Since the centre is
a Cartier curve in $P'$, the actual
restriction of its blowup is an
isomorphism $P\to P'$. In particular
it identifies the section $C$ with
the original curve $B$.

\par\medskip\noindent\textbf{Each ruling-fibre boundary becomes a transverse double curve.}
Consider one of the other image components
$Q'=\pi(Q)$. It does not contain $B$:
otherwise its generic first normal
divisor would meet every ruling fibre,
whereas $Q\cap T$ is a finite union
of fibres. More explicitly, writing
$\pi^*Q'=Q+aT$ and restricting to a
generic ruling fibre gives $a=0$.
Thus $\pi^*Q'=Q$.

Restrict a local defining section of
$Q'$ to $B$. Its pullback to $T$
vanishes exactly on the specified
reduced fibres belonging to $Q$.
At $b_j$ that section therefore
has a simple zero on $B$. Its
differential is nonzero on the
tangent line of $B$. Consequently
$Q'$ is smooth at $b_j$ and is
transverse to $B$. The independent
normal differential of $P'$ makes
$P'$ and $Q'$ transverse as well.

There is only one such other
component at each $b_j$. A second
would have the same fibre in
$T\cap Q$, contradicting the
given SNC intersections and
(C.RC.1). Any further image
component passing through $b_j$
would likewise have a boundary
component on the exceptional
fibre; hence none is missing
from this count. Thus there is
no triple point of the descended
support on $B$.

Choose coordinates with $P'$ and
$Q'$ the first two coordinate
planes and $B$ the stated
coordinate curve. Transversality
just proved permits this choice.
The full fibre is a unit times
$p_0^m q_0^{n_j}$; a local root
of that unit absorbs it into
$p_0$ without changing these
marked loci. Blow up $(p_0,r_0)$.
The exact chart $p_0=xy,r_0=x$
gives (C.RC.3). On $Q'$ this
ideal is the maximal ideal
of $b_j$, so $Q$ is precisely
its ordinary point blowup,
with exceptional fibre $F_j$.
On $P'$ the ideal is Cartier,
giving the already established
isomorphism. These are the
claimed original neighbourhood
maps.

\par\medskip\noindent\textbf{The full normal extension and the weighted line.}
The normal sequence of
$B\subset P'\subset U'$
is the first row of (C.RC.4).
The exceptional divisor of
the ordinary blowup is the
lines projectivization of
this very $V$, with
ambient normal $\OO(-1)$.
The strict parent intersects
it in the specified subline
$\PP_B(L)$. In particular,
with these line conventions,
\[
                 \OO_T(C)\simeq
                      \OO_{\PP_B(V)}(1)\otimes p^*M.
\tag{C.RC.10}
\]
The principal full fibre on
$U'$ is $mP'+\sum_Qm_QQ'$.
Restriction to $B$ gives
exactly the last isomorphism
of (C.RC.4), since each
$Q'$ meets $B$ simply at
its listed $b_j$. This
also follows by substituting
(C.RC.10) into (C.RC.8).
Taking degrees proves the
stated divisibility. The
line $M$ and its actual
power trivialization are
retained, including any
torsion information.

Restrict the actual vertical
canonical evaluation section
on $U'$ to $B$. Its line
isomorphism is
\[
 K_{U'}|_B\simeq
       M^{k_P}\otimes
             \OO_B\left(\sum_jk_{Q_j}b_j\right).
\tag{C.RC.11}
\]
Other components have a
nonvanishing defining section
on $B$ and contribute their
specified trivializations.
Adjunction for $B\subset P'$
then proves the formula for
$L$ in (C.RC.5). The
ordinary curve-blowup
Jacobian contributes $T$
once to $K_U$; the generic
centre lies only in $P'$.
Thus $k_T=k_P+1$.
Changing the canonical
generator by a power of $f'$
adds the full multiplicities
to the orders, and (C.RC.4)
shows that (C.RC.5) remains
the same line isomorphism.

In the full reduced case
$m=n_j=1$, so $M=\OO_B(-Z)$.
The derivative of
$f'=p_0q_0$ along the
normal bundle of $B$
is the quotient normal
map to $M$, followed by
its inclusion in $\OO_B$
with simple zero at $b_j$.
Away from the marked
points it is surjective.
This proves all maps in
(C.RC.6). For $m>1$ the
ordinary derivative of
$f'$ is zero along the
generic centre; the
quotient in (C.RC.4)
instead is the actual
normal map of $P'$,
and has already been
specified independently
of that derivative.

For the splitting assertion,
the extension group is the
actual space $H^1(B,L\otimes
M^{-1})$. Equations (C.RC.6)
and Serre duality identify
its dual with
\[
 H^0\left(B,\OO_B\left(
                   \sum_j(k_{Q_j}-k_P-2)b_j\right)\right).
\tag{C.RC.12}
\]
At a local maximum
$k_{Q_j}\le k_T=k_P+1$.
Every coefficient in
(C.RC.12) is at most
minus one. When $r>0$
the line has negative
degree and no section.
The extension group is
zero, which proves the
actual splitting with
the marked subline.
When $r=0$ the extension
is precisely the one
retained in UN64.

\par\medskip\noindent\textbf{Euler, triples, the converse and compact descent.}
The actual blowup gives
$S\setminus T\simeq S'\setminus B$,
whereas $T$ is a
projective-line bundle
of Euler number $2e(B)$.
Hence $e(S)-e(S')=e(B)$,
independently of the
full multiplicities.
At each $b_j$ the chart
(C.RC.3) has exactly one
triple point upstairs
and only a smooth double
curve downstairs. There
are no other changes,
which proves (C.RC.7).

Conversely, blowing up the
indicated transverse curve
uses exactly these charts.
The exceptional multiplicity
is the generic parent
multiplicity $m$, and
every other weight is
unchanged. The section
and ruling fibres are
exactly the intersections
in (C.RC.1), so all
the displayed maps and
markings are recovered.
As in UN64, the actual
local blowdowns glue to
the identity off $T$.
This proves the compact
version as well.
\end{proof}
